\subsection{Maximum \& Minimum Values of Sinusoidal Functions}
Similiar to polynomial function. 

The local maximum of a trigonometric function can be found by finding the critical numbers then testing regions to the 
left/right of these values, or using the second derivative. 

But please pay some of the attention to the difference, the intervals of increase/decrease in \textbf{Sinusoidal} functions are periodic. 

\begin{example}
    Determine the 'x' value which the local maximum values Occur for $$y = 3\cos^2 (x + \frac{\pi}{6}) - 2$$

    Solution:

    Find the derivative of $y$ first:
    \begin{align*}
        \frac{dy}{dx} &= 6\cos (x + \frac{\pi}{6}) \sin (x + \frac{\pi}{6}) * (-1) * 1\\
        &= -3 \left[2\cos(x + \frac{\pi}{6})\sin (x + \frac{\pi}{6})\right] \\
        &= -3\sin [2(x + \frac{\pi}{6})]\\
        &= -3 \sin (2x + \frac{\pi}{3})
    \end{align*}

    Then we solve for the critical points:
    \begin{align*}
        0 &= \sin (2x + \frac{\pi}{3})
    \end{align*}

    There are two possible cases:

    Case 1:
    \begin{align*}
        2x + \frac{\pi}{3} &= \arcsin 0\\
        2x + \frac{\pi}{3} &= 0\\
        2x &= -\frac{\pi}{3}\\
        x &= -\frac{\pi}{6}
    \end{align*}

    Case 2:
    \begin{align*}
        2(x + \frac{\pi}{6}) &= \pi - \arcsin 0\\
        2(x + \frac{\pi}{6}) &= \pi\\
        x &= \frac{\pi}{2} - \frac{\pi}{6} \\
        x &= \frac{\pi}{3}
    \end{align*}

    Because the period for the original function is $\pi$, so we actually have two more solutions.
    \begin{align*}
        x &= -\frac{\pi}{6} + \pi \\
         &= \frac{5\pi}{6}
    \end{align*}
    \begin{align*}
        x &= \frac{\pi}{3} + \pi \\
         &= \frac{4\pi}{3}
    \end{align*}

    Next, you can do the line test and determine the increase/decrease intervals. 
    Then, you can find out the maximum/minimum point. 
\end{example}

\subsubsection{Phyiscs Application}
You probably heard about \textit{Fourier Series} or \textit{Fourier Analysis}. 
There are lots of application of mathematics in the physics content. 
Two main types of question are discussed in the lecture.

\begin{center}
    \begin{itemize}
        \item Power Supplies question
        \item Simple Harmonic Motion
    \end{itemize}
\end{center}

\paragraph{Power Supplies} :
    The voltage of a power supply in an AC-DC couple circuit can be given by any Sinusoidal function. 
When solve for the maximum/minimum voltage for $V(t)$ function, apply the algorithm discussed in this section. 


\paragraph{Simple Harmonic Motion}

The horizontal position of a bob on a pendulum swings back and forth can be modeled by this equation:
\begin{displaymath}
    h(t) = A\cos(\frac{2\pi}{T}t)
\end{displaymath}
where:
\begin{center}
    \begin{itemize}
        \item $A = $ amplitude of the Simple Harmonic Motion
        \item $T = $ period of the Simple Harmonic Motion
    \end{itemize}
\end{center}

Period (aka "T") can be modeled by this equation:
\begin{displaymath}
    T = 2\pi\sqrt{\frac{l}{g}}
\end{displaymath}
where $l$ is the length of the pendulum string and $g$ is the acceleration due to gravity.

